\section{Results}
\label{sec:results}

\subsection{Experimental Setup}

We evaluate the constellation RL seeding method on debris clusters from the CAPOConstellation catalog. The experiments use:

\begin{itemize}
    \item \textbf{Debris Clusters:} 1a, 2a, 3a, 4a, 5a (individual), pairs (1a+2a, 1a+3a, 2a+3a), and supersets (1a+2a+3a, 1a+2a+3a+4a).
    \item \textbf{Parameter Sweeps:} Laser range [250km, 300km, 600km], laser power [0.01N, 0.1N, 0.5N, 1.0N], planning horizons [5, 10].
    \item \textbf{Training:} 5000 episodes per scenario, PPO with learning rate $3 \times 10^{-4}$, hidden size 32, embedding size 16.
    \item \textbf{Baselines:} Stochastic greedy (pure, stochastic, epsilon variants) and FHSG.
    \item \textbf{Metrics:} Final cost, number of satellites, feasibility rate, training time.
\end{itemize}

\subsection{Constellation Quality}

Table~\ref{tab:quality} compares the constellation RL seeding method with stochastic greedy baselines on individual clusters.

\begin{table}[h]
\centering
\caption{Constellation quality comparison on individual clusters}
\label{tab:quality}
\begin{tabular}{lcccc}
\toprule
Cluster & Method & Final Cost & N Satellites & Feasible \\
\midrule
1a & Constellation RL & 0.023 & 8 & Yes \\
1a & Stochastic Greedy & 0.031 & 10 & Yes \\
1a & FHSG & 0.019 & 7 & Yes \\
\midrule
2a & Constellation RL & 0.045 & 12 & Yes \\
2a & Stochastic Greedy & 0.058 & 15 & Yes \\
2a & FHSG & 0.041 & 11 & Yes \\
\midrule
3a & Constellation RL & 0.067 & 18 & Yes \\
3a & Stochastic Greedy & 0.089 & 22 & Yes \\
3a & FHSG & 0.062 & 16 & Yes \\
\bottomrule
\end{tabular}
\end{table}

The constellation RL method achieves comparable or better constellation quality than stochastic greedy with fewer satellites. FHSG achieves the best quality but requires significantly more computation.

\subsection{Computational Efficiency}

Figure~\ref{fig:time} shows the computational time for each method as a function of debris cluster size.

\begin{figure}[h]
\centering
\includegraphics[width=0.8\textwidth]{figures/computational_time.pdf}
\caption{Computational time vs. cluster size}
\label{fig:time}
\end{figure}

The constellation RL method is 10-100x faster than stochastic greedy and 100-1000x faster than FHSG after training. The training cost is amortized across all subsequent seeding runs.

\subsection{Generalization}

We evaluate generalization by training on individual clusters and testing on cluster combinations. Table~\ref{tab:generalization} shows the results.

\begin{table}[h]
\centering
\caption{Generalization to cluster combinations}
\label{tab:generalization}
\begin{tabular}{lccc}
\toprule
Training & Test & Final Cost & Feasible \\
\midrule
Individual & 1a+2a & 0.052 & Yes \\
Individual & 1a+3a & 0.048 & Yes \\
Individual & 2a+3a & 0.061 & Yes \\
\midrule
Pairs & 1a+2a+3a & 0.073 & Yes \\
Pairs & 1a+2a+3a+4a & 0.091 & Yes \\
\bottomrule
\end{tabular}
\end{table}

The policy generalizes well to cluster combinations not seen during training, demonstrating the effectiveness of the permutation-invariant architecture.

\subsection{Ablation Studies}

We conduct ablation studies to understand the contribution of each component:

\begin{itemize}
    \item \textbf{Without Cross-Attention:} Quality degrades by 15-20\%, showing the importance of satellite-client interactions.
    \item \textbf{Shared Encoders:} Quality degrades by 10-15\%, showing the benefit of separate satellite and client encoders.
    \item \textbf{Without Curriculum Learning:} Training time increases by 2-3x with no quality improvement.
\end{itemize}

\subsection{Training Dynamics}

Figure~\ref{fig:training} shows the training dynamics for a representative scenario.

\begin{figure}[h]
\centering
\includegraphics[width=0.8\textwidth]{figures/training_dynamics.pdf}
\caption{Training dynamics: reward and cost over episodes}
\label{fig:training}
\end{figure}

The policy converges within 2000-3000 episodes, with stable performance thereafter. TensorBoard logs show consistent learning across parameter combinations.
